\section{Experiments}
\label{sec:experiments}

We justify each component of the final system with leaderboard evidence. Many comparisons move more than one factor at once; we flag the confounds where they occur and treat this as a development trace rather than a single controlled ablation.

\subsection{Evaluation Protocol}
\label{sec:eval_protocol}

The task is scored by the macro averaged per quadrat F1 of \S\ref{sec:introduction} on a hidden test set of $2{,}105$ vegetation quadrats. Kaggle reports a public score on a fixed subset of quadrats and a private score on the complement; the two do not always rank configurations the same way (\S\ref{sec:infer_ablations}), so we report both wherever a configuration was submitted. Repeated near identical submissions put the run to run noise floor at $\approx 0.002$ Macro F1, and we treat smaller differences as not significant. We also track held out top-1/top-5 accuracy on a $10\%$ single plant validation split, but it correlates only weakly with quadrat Macro F1 ($r\approx0.4$; \S\ref{sec:val_lb}), so we use it to choose checkpoints, not to rank submissions.

\subsection{Baselines and Supervised Fine-tuning}

Off the shelf BioCLIP features go only so far: zero shot tile classification reaches ${\sim}0.12$ public Macro F1, a frozen backbone linear probe $0.242$, and a few shot support bank $0.307$. Neither zero shot matching nor a frozen backbone bridges the single plant to multi species shift, so we carry only the fine tuning line forward. Partially fine tuning BioCLIP~2.5 ViT-H/14 with the taxonomy aware head gives the decisive gains: the first strong model (experiment~010: shared trunk, five levels, $1.41$M images) reached $0.391$ public Macro F1, and our final model (experiment~i002: per head MLPs, three levels, $2.65$M images, ten epochs) reached $0.41165$ public / $0.38132$ private at single resolution inference, before calibration and ensembling. The comparison changes data, schedule, head topology, and head count at once, so we report i002 as the stronger configuration without isolating a single cause.

\subsection{Backbone Unfreezing}
\label{sec:blocks_ablation}

The one factor we swept within a single experiment is the number of unfrozen transformer blocks. Within experiment~010 ($1.41$M images) public Macro F1 rises from the head only baseline ($0.362$) to a peak at eight blocks ($0.391$) and declines at twelve, while the full fine tune collapses to $0.301$ public, trading the pretrained representation for single plant fit faster than the leaderboard rewards. Validation top-5 keeps rising with depth even as the leaderboard score peaks and falls, the decoupling of \S\ref{sec:val_lb}. Our final model (experiment~i002) repeated this on the larger $2.65$M manifest, but the optimum is shallower: public Macro F1 peaks at four blocks (Table~\ref{tab:blocks_i002}), which we adopt as the anchor, and the bottom rows show the inference steps that take it to the headline $0.418$.

\begin{table}[ht]
  \centering
  \small
  \begin{tabular}{lcc}
    \toprule
    i002 stage (2.65M manifest) & Best public F1 & Best private F1 \\
    \midrule
    head only (0 blocks)         & 0.394 & 0.388 \\
    \textbf{4 blocks (anchor)} & \textbf{0.412} & 0.389 \\
    8 blocks                     & 0.405 & 0.375 \\
    \midrule
    \multicolumn{3}{l}{\emph{inference enhancements on the 4-block checkpoint}}\\
    \quad + logit adj.\ ($\tau{=}0.25$, $T{=}0.03$) & 0.406 & \textbf{0.406} \\
    \quad + $224{+}336$ ensemble (anchor)            & \textbf{0.418} & 0.403 \\
    \bottomrule
  \end{tabular}
  \caption{i002 staged unfreeze sweep on the $2.65$M image manifest (ten epochs per stage). Unlike the 010 sweep, public Macro F1 peaks at four blocks, so the four block checkpoint was submitted. The bottom rows add the logit adjustment and resolution ensemble inference steps that produce the headline anchor.}
  \label{tab:blocks_i002}
\end{table}

\subsection{Tile Aggregation, Calibration, and Ensembling}
\label{sec:infer_ablations}

Of the three pooling rules of \S\ref{sec:adaptive_inference}, softmax mean wins: on the i002 four block checkpoint it beat max on both splits ($0.41165$ vs.\ $0.37509$ public), and plain mean was far weaker on a matched 010 checkpoint ($0.258$ against $0.340$). \label{sec:tile_aggregation} We then fixed the remaining knobs by leaderboard sweep. A joint sweep over $\tau$ and $T$ has a sharp optimum at $\tau{=}0.25$, $T{=}0.03$ ($0.40590$ public, falling to $0.36423$ at $T{=}0.05$ and $0.38440$ at $\tau{=}0.5$); at this $\tau$ the adaptive threshold also beats fixed top-$k$, which is second best at $k{=}2$ ($0.39878$, matching the two to three species per quadrat cardinality). These values were chosen by following the public leaderboard rather than on a held out split, so they are a reasonable operating point rather than a tuned optimum. Finally, averaging the same checkpoint at $224$ and $336$ pixels lifts public Macro F1 from $0.41165$ to $0.41826$ ($+0.0066$) and private from $0.38132$ to $0.40283$ ($+0.022$). The best public and best private operating points differ: logit adjusted single resolution scores $0.40590$ public but $0.40600$ private, the project's strongest private result, so we carry both columns through.

\subsection{Final Submission}
\label{sec:final_results}

The headline configuration scores \textbf{0.41826} public / $0.40283$ private Macro F1, the best public result of the project. Three architecture orthogonal pivots each rerank the anchor posterior with an additional signal, and none improves it (Table~\ref{tab:summary}). The triple backbone late fusion (cRT), a feature level concatenation of frozen BioCLIP~2.5, DINOv3, and ConvNeXt-V2-L features mixed with the anchor at $w{=}0.3$, was the worst ($\Delta{=}-0.036$): a pre-submission gate had already flagged the head as structurally divergent (top one agreement $0.094$), diversity going against signal. The genus/family co-occurrence rerank, a log prior built from the anchor's own aggregated supports, swapped $11\%$ of selections but barely moved the score, since a prior derived from the model's own logits carries little new information. The seasonal phenology prior, analysed in Section~\ref{sec:phenology}, was likewise near neutral. We interpret this pattern in \S\ref{sec:orthogonality}.

\begin{table}[ht]
  \centering
  \small
  \setlength{\tabcolsep}{4pt}
  \begin{tabular}{@{}p{0.58\linewidth}cc@{}}
    \toprule
    Config & Macro-F1 & $\Delta$ \\
    \midrule
    BC2.5 last 4 blocks, 224+336+LA, $T{=}0.03$ & \textbf{0.41826} & --- \\
    \hspace{1em}+ cRT triple, $w{=}0.3$ & 0.38195 & $-0.036$ \\
    \hspace{1em}+ Genus $\alpha_g{=}0.1$ & 0.41246 & $-0.006$ \\
    \hspace{1em}+ Pheno., $\beta{=}1.0$ & 0.41346 & $-0.005$ \\
    \bottomrule
  \end{tabular}
  \caption{All four submitted configurations (public Macro F1). The cRT pivot is strongly negative; the genus rerank and phenology prior move the anchor by only $-0.006$ and $-0.005$, small but consistent in sign. We read this as the BioCLIP~2.5 posterior already absorbing both signals through training distribution co-occurrence (\S\ref{sec:orthogonality}).}
  \label{tab:summary}
\end{table}
